\chapter{La reparación: códigos como numerales}
\label{cap:numerales}

El capítulo anterior terminó en un sitio incómodo: cuatro reparaciones razonables, y ninguna
funciona. Éste cuenta la que sí, y por qué —vista desde el final— era la única que podía funcionar.

\section{El diagnóstico, en una línea}

La inconsistencia no venía de un axioma equivocado sino de un \textbf{error de categoría}. El
símbolo objeto \ident{tcFn} —«el código del código»— es una operación sobre \emph{sintaxis},
declarada como función \emph{objeto}. Y una función objeto sólo puede depender de \textbf{valores}.

\begin{leccion}[La distinción que lo explica todo]
Un predicado \emph{es} un subconjunto: separa números, no maneras de escribirlos.
\ident{substfc} necesita un reconocedor \textbf{extensional} —¿pertenece este número al conjunto de
los códigos de fórmula?—, y eso un predicado lo puede dar. \ident{tcFn} necesita una distinción
\textbf{intensional} —¿qué sintaxis escribimos para denotar el número 9?—, y eso ningún predicado
lo puede dar, porque no está separado en los valores.
\end{leccion}

Con $\codigo{\varphi}$ escrito como \textbf{árbol} \ident{cons}, la sintaxis no se recupera del
valor: \ident{ax_L0_cons_def} identifica \ident{cons h t} con \ident{pair h (σt)}, así que en
$\mathbb{N}$ el mismo número es a la vez un par y un sucesor, y \ident{tcFn} debía devolver dos
códigos distintos para él. Con $\codigo{\varphi}$ escrito como \textbf{numeral}, la ambigüedad
desaparece por una razón que conviene decir despacio: \textbf{el numeral es canónico}. Sólo hay una
manera de escribir $\sigma^n 0$.

\begin{refutado}[Y no era «indemostrable»: era falso]
La lectura sintáctica de \ident{tcFn} sobre \ident{cons} con argumentos abstractos no es una ecuación
que faltase probar. Bajo la lectura numeral obliga a que un código de cabeza $\codigo{\sigma}$ sea
igual a uno de cabeza $\codigo{::}$. Medido en \modulo{sondeos/PilotoRastreada.lean}.
\end{refutado}

\section{La aritmética sale sin división}

El obstáculo práctico de la vía numeral es el polinomio de Cantor: \ident{cons} se define como
$\mathrm{div2}$ de él, y razonar sobre divisibilidad dentro de la teoría objeto es caro. La salida
es no razonar sobre ella en absoluto. Se define el valor del código con \textbf{números
triangulares}, de modo que la mitad ya está tomada.

\begin{matematica}[Números triangulares y valor de un \ident{cons}]
\[
T(0)=0,\qquad T(n{+}1)=T(n)+(n{+}1),
\qquad\text{luego}\qquad 2\,T(n)=n(n{+}1).
\]
\[
\mathrm{consN}(a,b) \;=\; T\bigl(a+(b{+}1)\bigr) + (b{+}1),
\qquad
2\cdot\mathrm{consN}(a,b) \;=\; \underbrace{(a{+}b{+}1)(a{+}b{+}2) + 2(b{+}1)}_{\text{polinomio de Cantor de }\langle a,b\rangle}.
\]
La segunda identidad es una igualdad de \ident{Nat}, demostrable por reescritura. No hay división
en ninguna parte.
\end{matematica}

\leanfrag{triN}
\leanfrag{consN}
\leanfrag{two_mul_consN}

\begin{leccion}[Dónde está el truco]
No se ha evitado la división eligiendo mejor las tácticas: se ha evitado \emph{eligiendo la
definición}. \ident{consN} se define ya multiplicado por dos, y por eso engancha directamente en la
forma $2m$ exacta que pide \ident{prf_div2_numeral}. Cuando una prueba obliga a razonar sobre una
operación cara, la primera pregunta no es cómo probarlo, sino si la definición se puede escribir
del otro lado.
\end{leccion}

\section{De la identidad en \ident{Nat} a la evaluación provable}

La identidad anterior vive en la metateoría. Para que sirva hay que subirla al cálculo: la teoría
objeto tiene que \emph{probar} que el árbol de código y el numeral de su valor son el mismo número.

\leanfrag{prf_cons_eval}

Con ese peldaño, el espejo meta de \ident{formCode} —la función \ident{codeNat}, que calcula en
\ident{Nat} lo que \ident{formCode} construye como término— cierra el puente por meta-recursión sobre
la estructura de la fórmula.

\leanfrag{codeNat}

\begin{teorema}[El código de una fórmula \emph{es} un numeral]
\label{teo:formcode-numeral}
Para toda fórmula $\varphi$, la teoría prueba
$\codigo{\varphi} \;=\; \mathrm{numeral}\bigl(\mathrm{codeNat}(\varphi)\bigr)$.
\end{teorema}

\leanfrag{prf_formCode_numeral}

Obsérvese que la prueba es una \textbf{meta-recursión}: un caso por constructor de \ident{Formula},
cada uno ensamblado con \ident{prf_cons_eval_of}. No hay inducción de la teoría objeto por ninguna
parte, y por eso el footprint se queda en la base del kernel.

\section{El puente, y qué cambia en los enunciados}

Del lado del cálculo $\derives$, el mismo hecho es la hipótesis que el piloto del lema diagonal
daba por supuesta, y que aquí queda \textbf{descargada}:

\leanfrag{hFN}

\begin{muro}[Lo que sí cambia]
Refundar por la vía numeral \textbf{cambia enunciados, no sólo pruebas}: el código estático que
aparece dentro de los teoremas pasa de \ident{termCode (formCode φ)} a
\ident{termCode (numeral (codeNat φ))}. Cada consumidor hay que revisarlo, o componerlo con
\ident{provCode_transfer}, que puentea las dos representaciones en un solo paso de Leibniz.
\end{muro}

\section{Cómo se decidió: pilotar antes de ejecutar}

Aquí está, para mi gusto, la lección de ingeniería del capítulo. La reparación exigía construir
unos veinte lemas nuevos. Antes de escribir ninguno se hizo lo contrario de lo intuitivo:
\textbf{se asumió el resultado costoso como axioma de Lean} y se comprobó que, con él, la cadena
cerraba de verdad.

\leanfrag{piloto_codeN_axioma}
\leanfrag{piloto_hFN_axioma}

Con esos dos axiomas puestos, el piloto reconstruyó el lema diagonal entero y comprobó que el punto
fijo salía. Sólo entonces se invirtió el trabajo en probarlos — y el sondeo gemelo descargó \ident{hFN}
sin axioma ninguno:

\leanfrag{descarga_hFN}

\begin{leccion}[Un axioma temporal es un instrumento de medida]
Poner un \ident{axiom} de Lean no es hacer trampa si se hace \emph{para preguntar}. Contesta la única
pregunta que importa antes de invertir: \emph{si tuviera esto, ¿cerraría?} Si la respuesta es no, se
han ahorrado veinte lemas; si es sí, se sabe exactamente dónde van a enchufarse. Lo que sería hacer
trampa es olvidarse de quitarlo — y contra eso está el inventario de \doc{AXIOMS.md} y la auditoría
por \printaxioms.
\end{leccion}

\section{El coste}

\begin{center}\small
\begin{tabular}{ll}
\toprule
Axiomas retirados & \ident{ax_tc_cons} \\
Axiomas nuevos & ninguno \\
Balance & $-1$ \\
Módulos afectados & 31, a cuarentena (capítulo siguiente) \\
\bottomrule
\end{tabular}
\end{center}

\noindent Es un resultado poco frecuente: la reparación de una inconsistencia que \textbf{reduce} la
superficie axiomática en vez de ampliarla. La razón es que no se añadió nada para tapar el agujero;
se cambió la \emph{representación} para que el agujero no existiera.

\avisoreparacion
